% !TeX program = xelatex
% Self-contained source. Compile twice with XeLaTeX.
\documentclass[11pt,a4paper]{article}
\usepackage{fontspec}
\usepackage[russian]{babel}
\IfFontExistsTF{Times New Roman}{\setmainfont{Times New Roman}}{\IfFontExistsTF{Liberation Serif}{\setmainfont{Liberation Serif}}{\setmainfont{CMU Serif}}}
\IfFontExistsTF{Arial}{\setsansfont{Arial}}{\setsansfont{DejaVu Sans}}
\IfFontExistsTF{Consolas}{\setmonofont{Consolas}[Scale=0.88]}{\setmonofont{DejaVu Sans Mono}[Scale=0.82]}
\usepackage[margin=22mm,headheight=14pt]{geometry}
\usepackage{amsmath,booktabs,tabularx,longtable,graphicx,xcolor,hyperref}
\definecolor{navy}{HTML}{173A56}
\definecolor{pale}{HTML}{EFF4F7}
\hypersetup{colorlinks=true,linkcolor=navy,urlcolor=navy,citecolor=navy,
  pdftitle={DriveMem: отчёт об экспериментах и проверенных гипотезах},
  pdfauthor={Исследовательский проект DriveMem}}
\setlength{\parindent}{1.2em}
\setlength{\parskip}{3pt}
\setlength{\emergencystretch}{2em}
\renewcommand{\arraystretch}{1.16}
\linespread{1.055}
\Urlmuskip=0mu plus 2mu
\def\UrlBreaks{\do\/\do\_\do\-\do\.}
\newcommand{\code}[1]{\texttt{\detokenize{#1}}}
\newcommand{\gain}{\Delta Q}
\newcommand{\resultbox}[1]{\par\medskip\noindent\fcolorbox{navy}{pale}{\parbox{\dimexpr\linewidth-2\fboxsep-2\fboxrule\relax}{#1}}\par\medskip}
\DeclareMathOperator*{\argmax}{arg\,max}
\setcounter{tocdepth}{1}

\begin{document}
\begin{titlepage}
\noindent{\sffamily\color{navy}\Large ИССЛЕДОВАТЕЛЬСКИЙ ОТЧЁТ}\par
\vspace{18mm}
{\sffamily\bfseries\color{navy}\Huge DriveMem}\par
\vspace{5mm}
{\sffamily\LARGE Эпизодическая память для\par переоценки траекторий DriveSuprim}\par
\vspace{10mm}
{\Large Реализованные методы, проверенные гипотезы,\par положительные и отрицательные результаты}\par
\vspace{12mm}
\noindent\rule{\linewidth}{0.8pt}
\vspace{4mm}
\noindent\begin{tabularx}{\linewidth}{@{}lX@{}}
Дата среза & 7 ноября 2026 года \\
Экспериментальная основа & DriveSuprim R34, NAVSIM v2; отдельная диагностика v1 \\
Финальная абляция & \code{component_decoder_functional_v1} \\
Тип документа & Ретроспективный технический и научный отчёт \\
\end{tabularx}
\vfill
\resultbox{\textbf{Основной итог.} В проверенных пулах есть существенный oracle-резерв выбора. Построен воспроизводимый стенд памяти и контроля вмешательств. Однако устойчивое преимущество реализованных memory-методов над baseline и сопоставимыми no-memory моделями не подтверждено. В последней абляции заранее заданное правило выбора сохранило \code{primary=baseline}.}
\vfill
\noindent\small Отчёт подготовлен по коду проекта, исходному плану статьи и серверным выводам, переданным пользователем. Серверные эксперименты запускал пользователь в контейнере \code{egorov_drivesuprim}; код, проверки и разбор результатов разрабатывались в ходе совместной работы. Прямое повторное исполнение серверных экспериментов при подготовке этого текста не выполнялось.
\end{titlepage}

\tableofcontents
\clearpage

\section{Краткий итог и границы отчёта}
\subsection{Что удалось получить}
Работа прошла путь от проверки воспроизведения базового планировщика до нескольких вариантов переноса результатов прошлых действий, прямой регрессии полезности, компонентной переоценки и candidate-conditioned decoder-признаков. Исходный checkpoint не дообучался: изменялись внешние признаки, модели переоценки и правила выбора среди готовых кандидатов.

К положительным результатам относятся:
\begin{enumerate}
\item Получен результат DriveSuprim R34 на NAVSIM v2: $83{,}0677$ EPDMS, совпадающий с ориентиром $83{,}1$ после округления. Для v1 получено частичное, но не полное воспроизведение.
\item Показано наличие лучших траекторий в выдаче модели: например, на 64 сценах из 32 лог-сегментов oracle по 256 кандидатам дал $+5{,}7265$ п.п. локального PDM.
\item Реализованы разделение по сессиям, отдельные роли данных, raw/filtered labels, неизменяемые manifests и hashes, возобновляемые cache/export/label/fit/evaluate этапы.
\item Удалось уменьшить стоимость подготовки пилота: вместо предложенной передачи до $449$ ГБ использованы бюджетные архивы, а расширение M/F выполнено загрузкой одного дополнительного архива $8{,}43$ ГБ.
\item Проведены сравнения с no-memory, random-compatible и shuffled-outcome контролями. Показано, что дополнительные метки и более богатые признаки сами по себе не обеспечивают перенос.
\item После фиксации вычислительного режима оба текущих forward совпали точно во всех 152 сценах аудита. Сохранены baseline, top16, vocabulary и соответствие decoder latent штатным heads.
\end{enumerate}

\subsection{Что не удалось подтвердить}
Ни одна из последних двух компонентных постановок не прошла собственное правило выбора полезной основной модели на F-CV: primary остался baseline. Decoder-memory дала $+4{,}0653$ п.п. на первой development-когорте, но весь положительный эффект обеспечила одна уже известная сцена. На второй когорте guarded-вариант не вмешивался, а component-вариант не превзошёл no-memory контроль.

Не получены доказательства калиброванного контроля риска, пользы растущего во времени банка, переноса на второй planner, улучшения полного EPDMS с памятью или closed-loop эффекта. Эти пункты остаются задачами исходного плана, а не результатами выполненной работы.

\subsection{Как трактовать достоверность чисел}
Основные серверные числа в этом документе взяты из переданных JSON/CSV-выводов, а не из независимого доступа к серверу. В комплект включены извлечённые JSON и manifest локально доступных источников с SHA256. Наличие \code{status=complete} в отчёте модели доказывает завершение соответствующего программного этапа, но не заменяет проверку научной гипотезы.

В тексте различаются: (а) непосредственно сообщённые результаты; (б) проверенная по коду конструкция метода; (в) интерпретации и рекомендации. Отрицательный результат небольшого пилота не является доказательством бесполезности любой эпизодической памяти.

\section{Исходный замысел статьи и фактически проверенная постановка}
\subsection{Цель исходного плана}
Исходный шестимесячный план задавал тему \emph{Similarity Is Not Utility: Risk-Controlled Episodic Recall for Continual Motion Planning}. Предлагалось сохранять прошлые эпизоды, извлекать релевантный опыт и оценивать пользу вмешательства до изменения решения планировщика. Предусматривались покомпонентные эффекты, доверительные границы, оценка риска вреда, независимая калибровка, рост памяти, новые маршруты и два планировщика~\cite{plan}.

Для frozen planner $\pi_0$ обозначим множество кандидатов через $C(s)$, baseline через $\tau_0(s)$, решение после обращения к памяти через $\tau_M(s)$. Планируемая целевая величина:
\begin{equation}
 \Delta_k(s,M_s)=U_k(\tau_M(s)\mid s)-U_k(\tau_0(s)\mid s).
\end{equation}
Критерий допуска из плана должен был опираться на нижние границы пользы критических компонентов и допустимый риск вреда. Реализованные в пилоте point-estimate gates не являются такими доверительными границами.

\subsection{Почему начали с ограниченного пилота}
До обучения памяти требовалось проверить более простую предпосылку: есть ли в готовом пуле альтернативы лучше штатного выбора. Иначе memory-reranker, не генерирующий новых траекторий, не сможет существенно улучшить результат. Затем проверялись дешёвые внешние модели с frozen представлением и уже имеющимися изображениями.

Реальная экспериментальная постановка уже исходного плана: фиксированный enriched банк, один planner, небольшие archive-budgeted выборки, offline local PDM, повторно просмотренный development. Визуальное представление исходной памяти относится к текущему входному изображению; ego-history присутствует, но полноценная модель временного эпизода не реализована. Точный поиск в небольшом банке использован вместо FAISS. Эти упрощения позволяют проверить отдельные механизмы, но ограничивают заявления об episodic и continual planning.

\section{Воспроизведение DriveSuprim и выбор основной платформы}
\subsection{NAVSIM v2}
В документации проекта зафиксирован EPDMS $83{,}0676746838$ для опубликованного R34 checkpoint. Результат соответствует $83{,}1$ после округления~\cite{repro}. Это положительный исход этапа reproduction, но совпадение одной агрегированной цифры само по себе не доказывает идентичность всех деталей авторского протокола.

\subsection{NAVSIM v1}
Локальная адаптация v1 дала PDMS $86{,}7300284806$ против ориентира $89{,}9$, то есть недостаёт примерно $3{,}1700$ п.п. Проверялись различия протоколов выбора и TTC. Две выдачи отличались в $5852/12146$ сценах, но сравнение на одном v1 scorer/cache дало только $86{,}7300\rightarrow86{,}7467$ ($+0{,}0167$ п.п.). HumanAgent дал $94{,}5455$ против ориентира $94{,}8$~\cite{repro}.

Отдельный контролируемый тест изменял только окно TTC при неизменных cache, simulated states и остальных компонентах. На 12146 сценах сокращение ego-окна с 4,0 до 3,1 с дало TTC $90{,}8859\rightarrow96{,}5009\%$ и PDMS $86{,}7300\rightarrow88{,}9221$ ($+2{,}1921$ п.п.). Это локализует значимую часть чувствительности протокола, но не воспроизводит авторский результат: TTC короткого варианта выше ориентира $93{,}6\%$, а PDMS всё ещё ниже $89{,}9$. Пересчёт старой агрегации по другому v2 CSV дал $88{,}8797$; это арифметический контраст разных прогонов, а не тот же причинный TTC-тест. Разбор EP показал, что основная разница его обычных средних приходится на сцены с нулевым multiplicative gate, а не доказывает ухудшение движения вперёд.\par

Причина полного расхождения с v1-таблицей не установлена. Иной checkpoint или авторские настройки остаются гипотезами. Проверки не привели к подмене checkpoint или произвольной поправке результата. Основные DriveMem-эксперименты были продолжены на v2, чтобы не смешивать нерешённый reproduction-вопрос и проверку новой памяти.

\section{Что было реализовано в экспериментальном стенде}
\subsection{Контейнер, данные и этапы}
Серверная рабочая среда: \code{egorov_drivesuprim}, репозиторий \path{/workspace/DriveSuprim}, эксперименты \path{/workspace/exp_v2}. Доступная GPU-конфигурация в переданных логах: GeForce GTX 1080, PyTorch $2.0.1+\mathrm{cu118}$, Python $3.9.23$. Начальный \path{/workspace/dataset} оказался read-only и содержал test-часть без trainval.

Были добавлены inventory и диагностика путей, подготовка navtrain metadata, запись в отдельный writable data root, отбор нужных изображений по frozen manifests, cache-status, экспорт candidates/features, label scoring, fit, evaluation и отчёты. Каждый дорогостоящий этап допускает возобновление. Готовые результаты проверяются по hashes, а не молча пересчитываются новым кодом.

\subsection{Подготовка данных с ограниченным бюджетом}
Инвентаризация сообщила 103288 доступных сцен, 162 исходные session groups и 159 eligible sessions. Первоначально выбраны 352 сцены в 64 сессиях. Для них cache оказался полным: $352/352$ корректных файлов. Но image plan требовал 11264 изображений и до $449{,}007$ ГБ сжатых архивов. Пользователь запросил существенно меньший объём.

Вместо этого был подготовлен \code{budget20gb_v1}: 136 сцен, 17 сессий и четыре исходных camera shards с суммарным сжатым размером около $17{,}148$ ГБ. Архивная доступность и стоимость влияли на отбор, поэтому это технический пилот, а не репрезентативный benchmark. Дополнительные четыре M/F-сессии позднее получены из shard 174 за $8{,}4276$ ГБ. Бюджет учитывает архивные байты; сетевые повторы и служебные накладные расходы не следует считать нулевыми.

\par
\begin{table}[htbp]\centering\small
\caption{Роли и размеры исходного бюджетного пилота. C/T outcomes оставались закрытыми.}
\begin{tabular}{@{}lrrl@{}}\toprule
Роль & Сессии & Сцены & Назначение\\\midrule
M & 6 & 48 & Донорский банк\\
F & 4 & 32 & Обучение и выбор параметров\\
V1 & 3 & 24 & Исходный development\\
C & 2 & 16 & Резерв калибровки\\
T & 2 & 16 & Резерв module holdout\\\midrule
Всего & 17 & 136 & \\\bottomrule
\end{tabular}
\end{table}
\par

Отдельный \code{validation_cached4_v1} содержит 16 сцен из четырёх новых относительно исходного пилота сессий. После первого просмотра и последующего выбора гипотез он является V2-development, несмотря на историческое слово validation в имени. После M/F-расширения обучающие данные составили M64/8 сессий и F48/6 сессий. Размеры V1/V2 не изменились.

\subsection{Изоляция ролей и происхождение supervision}
Session ID задаётся датой/временем и vehicle без двух конечных суффиксов log-segment. Разные сегменты одной session не разделяются между M/F/V/C/T. Интервал в 10 секунд уменьшает близость сцен, но не делает их статистически независимыми.

Банк \emph{simulator-enriched}: для прошлой сцены оценивались альтернативные действия, а не только реально выбранная траектория. No-memory модели получали те же M/F supervised labels. Таким образом, различие состоит в retrieval при формировании признаков, а не просто в дополнительном доступе к меткам. Равенство labels и процедуры обучения не означает полного равенства числа признаков или параметров.

Текущие query outcomes запрещены в retrieval и признаках. Решения методов сохраняются до lookup оценок соответствующих query-кандидатов. Исключается вся собственная M-сессия query, а не только совпадающий token; F не становится банком доноров. В shuffle-контроле исключение своей сессии выполняется до перестановки labels, чтобы собственный исход не переехал к другому донору.

\subsection{Неизменяемые артефакты и программные проверки}
Реализованы identities с code/checkpoint/config/data/runtime hashes, индексы экспорта, отдельные experiment roots и атомарная запись по сценам. Старые labels/features не заменяются при изменении формулы; эксперимент получает новый root. Проверялись отсутствие query-label consumption при принятии решений, собственных доноров, несогласованной нормировки folds, нарушения baseline fallback, перепутанных candidate IDs и ошибочного resume.

В конце разработки функционального режима локально прошли 35 decoder-related Python-тестов, включая проверки prior-audit provenance, и launcher-регрессии. Это свидетельство корректности проверенных программных условий, а не повторный GPU benchmark и не независимая научная валидация метода.

\section{Метрики, вред и статистика}
\subsection{Local PDM не равен полному EPDMS}
Все memory/oracle результаты ниже, если не сказано иное, относятся к \textbf{локальному simulated PDM без two-frame extended comfort}. Они не являются приростом EPDMS из полного reproduction-прогона. В используемой raw локальной конфигурации скалярная оценка имеет multiplicative gates и взвешенную часть. В частности, после исключения двухкадрового комфорта веса progress/TTC/lane keeping/history comfort равны $5/5/2/2$:
\begin{equation}
 Q_{\mathrm{local}}=G\,\frac{5\,EP+5\,TTC+2\,LK+2\,HC}{14}.
\end{equation}
Фактическое $Q$ для анализа берётся из сохранённого scorer output. Формула не используется для скрытого восстановления скаляра из уже изменённых human-filter компонентов. Нормировка progress зависит от стандартного механизма оценки; нельзя без проверки смешивать оценки из разных наборов proposals.

\subsection{Относительная польза и определение harm}
Для сравнения с baseline:
\begin{align}
 g_i &=100\,[Q_i(\tau_i)-Q_i(\tau_{0,i})],\\
 H_i &=\mathbf{1}\{\exists k\in\mathcal K_{\mathrm{safety}}:
 U_{ik}(\tau_i)+10^{-8}<U_{ik}(\tau_{0,i})\},\\
 \mathcal K_{\mathrm{safety}}&=\{NC,DAC,DDC,TLC,TTC\}.
\end{align}
Здесь $g_i$ измеряется в процентных пунктах на шкале $100Q$. Progress, comfort и lane keeping не включены в этот конкретный $H$, хотя сообщаются отдельно. Поэтому $H=0$ не исключает падения общего PDM. \code{quality_degraded} именно фиксирует снижение $Q$.

\emph{Coverage} есть доля сцен, в которых изменился candidate ID. Смена траектории с тем же итоговым баллом учитывается как вмешательство с нулевой пользой. Нулевой harm при нулевом coverage означает отсутствие вмешательств, а не полезную память. Абсолютное нарушение и относительный harm различаются: одинаково плохие baseline и новое действие могут дать $H=0$.

\subsection{Единица усреднения}
Основной агрегат в memory-отчётах:
\begin{equation}
 \overline g_{\mathrm{macro}}=\frac{1}{S}\sum_{j=1}^{S}
 \frac{1}{n_j}\sum_{i\in\mathrm{session}\ j}g_i.
\end{equation}
На V1 и V2 число сцен внутри соответствующей когорты одинаково по сессиям, поэтому macro и micro совпали. Bootstrap выполнялся по session blocks; интервалы при трёх или четырёх сессиях имеют ограниченную информативность. Leave-one-session-out sensitivity без переобучения является разбором зависимости эффекта от одной поездки, а не новым cross-validation экспериментом.

\section{Oracle-эксперименты: существует ли резерв выбора}
\subsection{Первый диагностический лог}
На восьми сценах одного лог-сегмента были оценены 256 кандидатов на сцену. Baseline local PDM $0{,}651683$, oracle $0{,}988179$, выигрыш $33{,}6496$ п.п. Четыре сцены уже имели балл 1; в четырёх других oracle существенно улучшал NC и TTC. Например, в одной сцене NC менялся $0{,}5\rightarrow1$, TTC $0\rightarrow1$, а локальная оценка $0{,}2609\rightarrow0{,}9054$.

Лучшие кандидаты имели model-score ранги 4, 51, 14 и 11 из 256. Для baseline head-прогнозы NC/TTC могли быть выше, чем для фактически более качественной альтернативы. Например, на сцене с рангом oracle 51 прогнозы NC/TTC были $0{,}9629/0{,}9529$ у baseline против $0{,}9210/0{,}8636$ у oracle, тогда как измеренные NC/TTC улучшались. Это мотивирует переоценку, но не доказывает, что похожий донор даёт нужную коррекцию.

\subsection{Расширение на 32 лог-сегмента}
После исключения первого пилотного лога выбраны 64 сцены: по две на каждый из 32 log-segments, seed 20261027, temporal gap 10 секунд. Получены:
\par
\begin{table}[htbp]\centering\small
\caption{Oracle по 256 кандидатам на 64 сценах; baseline/oracle на шкале $100Q$.}
\begin{tabular}{@{}lr@{}}\toprule
Baseline local PDM & 90,2431\\
Oracle local PDM & 95,9696\\
Средний выигрыш, п.п. & 5,7265\\
Улучшенные сцены & 32/64\\
Log-segments с улучшением & 23/32\\
Медианный выигрыш сцены, п.п. & 0,1461\\
Максимальный выигрыш сцены, п.п. & 97,5562\\
95\% log-cluster bootstrap CI, п.п. & [1,4896; 11,1356]\\\bottomrule
\end{tabular}
\end{table}
\par

Для model top-$k$ резерв составил:
\begin{center}\small
\begin{tabular}{@{}lrrrrr@{}}\toprule
$k$ & 8 & 16 & 32 & 64 & 256\\\midrule
Gain, п.п. & 4,1734 & 4,8106 & 5,2631 & 5,5112 & 5,7265\\\bottomrule
\end{tabular}
\end{center}
Это обосновало выбор top16 как дешёвого пилотного пространства действий. Однако последующий V2 показал, что top8 не всегда достаточен: там oracle top8 дал лишь $0{,}8451$, а top16 $7{,}2464$ п.п.

Указанный ранний CI построен по log-segments. По группировке из локального плана 32 сегмента соответствуют примерно 21 session prefix; это более грубые зависимые блоки. В отчёте не подменяется исходный CI невычисленным session-CI. В новых memory-экспериментах единицей группировки уже служит session.

\subsection{Проверка human penalty filter}
Реконструкция $Q$ из отображаемых компонентов обнаружила расхождения в 758 из 16384 candidate rows, затрагивающие четыре сцены и три baseline. Диагностика показала замену lane keeping $0\rightarrow1$ в filtered представлении при неизменном сохранённом скаляре; типичная разница составляла $1/7\approx0{,}142857$.

Выполнена повторная raw/filtered оценка выбранных baseline/oracle во всех 64 сценах: 96 уникальных candidate evaluations. Средний выигрыш до и после фильтра сохранился $5{,}7265$ п.п.; baseline score changes, oracle score changes и safety disagreements равны нулю; выбранные oracle не ухудшали пять учитываемых safety-компонентов. Это проверка выбранных решений, \textbf{не} повторная raw-оптимизация всех 256 альтернатив. В дальнейшем raw/filtered состояния и snapshots сохранялись явно.

\subsection{Oracle на двух development-когортах}
\par
\begin{table}[htbp]\centering\small
\caption{Action-oracle внутри top16; baseline/oracle на шкале $100Q$.}
\begin{tabular}{@{}lrrrr@{}}\toprule
Когорта & Сцены/сессии & Baseline & Oracle & Gain, п.п.\\\midrule
V1 & 24/3 & 88,7651 & 95,4101 & 6,6449\\
V2 & 16/4 & 85,6927 & 92,9391 & 7,2464\\\bottomrule
\end{tabular}
\end{table}
\par
На V1 улучшились 11 сцен, session-bootstrap CI $[0{,}4131;14{,}5425]$ п.п.; на V2 улучшились 8 сцен, CI $[0{,}7717;19{,}5547]$. На этих наборах raw, filtered и соответствующие safe-oracle итоги совпали. Oracle использует исходы текущего запроса, недоступные методу при выборе; это верхний диагностический ориентир, а не deployable результат и не donor supervision.

\section{Первая память: перенос ошибок штатных heads}
\subsection{Идея и реализация}
Исходная гипотеза: в похожей сцене ошибки прогнозных heads будут похожими, поэтому их можно исправить переносом residuals. Для прошлого действия $a$ в эпизоде $e$:
\begin{equation}
 r_{e,k}(a)=y_{e,k}(a)-p_{e,k}(a),
\end{equation}
где $p$ --- sigmoid-выход головы, а $y$ --- измеренная компонента. Для query выбирались до пяти command-compatible M-доноров по frozen scene embedding и ego velocity. Действия сопоставлялись геометрическим ядром; чужая траектория непосредственно не копировалась.

Scene embedding получался усреднением frozen downscale-layer output. Были реализованы similarity memory, threshold по расстоянию, random-compatible donor и ridge без retrieval. Начальная разметка: baseline, model top8 и восемь случайных альтернатив, с устранением дубликатов.

\subsection{Результат первого пилота}
На V1 similarity memory и distance-threshold вариант дали нулевой gain при двух вмешательствах из 24. Random memory дал $+0{,}0286$ п.п. при трёх вмешательствах. Начальный no-memory ridge дал $-0{,}2614$ п.п., 16 вмешательств, один safety-harm и пять quality drops~\cite{initial}.

Причина нулевого выигрыша не сводилась к отсутствию доноров: совместимые доноры находились во всех 24 сценах, извлекались 45 из 48 M-эпизодов. Иными словами, найденный похожий опыт ещё не давал полезного изменения выбора. Это важный отрицательный результат, но малый размер и конкретное представление не позволяют утверждать универсальную несостоятельность similarity retrieval.

\section{Action transfer: геометрия действий и baseline contrast}
\subsection{Что изменили}
Следующая семья B0--B3 проверяла абсолютный и контрастный перенос, исходную и локальную ширину ядра. Contrast вычитает коррекцию baseline из коррекции альтернативы, чтобы baseline был опорной точкой. Масштаб $h_0$ определялся на M; $h$ выбирался по F среди локальных вариантов. Дополнены F top16 labels; сравнивались full-embedding ridge, scene-only residual и три random seeds. Все новые решения выбирались из одного top16~\cite{action}.

\subsection{Положительный результат и его концентрация}
B3 был выбран по F. На V1 он дал $+4{,}2356$ п.п., четыре вмешательства и нулевой наблюдаемый safety-harm. B1 дал $+4{,}5524$ п.п., no-memory ridge $+1{,}3234$. Однако random seeds дали $+4{,}0506$, $-0{,}0107$ и $+4{,}0893$ п.п. Это показывает значимую роль конкретного набора доноров и не даёт уверенно приписать выигрыш одному similarity retrieval.

Весь положительный gain B3 находился в сессии \code{2021.10.11.07.47.13_veh-50}; после её исключения средний эффект становился нулевым. Крупнейшее вмешательство --- \code{f27a618a5c7458ee}, $Q:0\rightarrow1$. Отдельный покомпонентный разбор был нужен, чтобы не объявить этот скачок исправлением столкновения без оснований.

\subsection{Проверка на V2}
На cached4 B3 дал $-0{,}9579$ п.п., четыре вмешательства, два quality drops и ни одной улучшенной сцены. Safety-harm при этом был нулевым. B0/B1/B2 дали ноль; no-memory ridge $-0{,}8742$ п.п. Один random seed имел небольшой плюс $+0{,}0805$, два остальных минус~\cite{cached}. Следовательно, хороший средний результат V1 не перенёсся на другую просмотренную development-когорту.

\section{Проверки причин неудачи residual transfer}
\subsection{Хватает ли donor labels в рабочем диапазоне действий}
На M все 384 слота рангов 1--8 были размечены, но из 384 слотов рангов 9--16 имелись только 15 labels ($3{,}91\%$). Большая часть старых случайных labels находилась за top16. Геометрически близкий action из top16 мог не иметь подходящего размеченного аналога~\cite{support}.

Доразмечены 369 отсутствующих M top16 действий; старые random labels сохранены. Общее число старых M labels стало 1137, F labels --- 756. Это отдельный sidecar, а не изменение исходных frozen артефактов.

\subsection{Ограничен ли выбор gate или предложениями доноров}
Выполнялись ретроспективные gate-oracle и family-oracle проверки. Gate-oracle может идеально решить, когда оставить baseline, но не может породить альтернативу, отсутствующую в предложениях. На V2 gate-oracle убирал отрицательный выигрыш B3 до нуля; all-family oracle проверенных исходных donor policies давал лишь $+0{,}5887$ п.п. при action-oracle $+7{,}2464$.

После полного M top16 completed-absolute all-family oracle достиг $+0{,}7327$ п.п. (около $10{,}1\%$ action-oracle резерва), retrieved-family oracle оставался около $+0{,}0884$. Крупный случай с action-oracle gain $+99{,}8080$ п.п. не исправлялся~\cite{capacity}. Эти capacity оценки относятся только к перебранным семействам, не ко всем возможным сочетаниям доноров; candidate-policy пары не являются независимыми наблюдениями.

\subsection{Не блокирует ли improvement сама формула score}
Anchored score-ceiling проверка не обнаружила запрета достижимости оптимального top16 action для восьми oracle-positive сцен V2: оптимум был формально достижим. Но oracle-head пробники выявили ограничение contrast-target: на V2 absolute-варианты при истинных query-компонентах дали около $+7{,}0589$ и $+7{,}0668$ п.п., contrast --- только $+0{,}8209$ и $+0{,}8288$~\cite{ceiling}.

Алгебраически при $\beta=1$, истинном residual $r(a)=y(a)-p(a)$ и contrast:
\begin{equation}
 p_{\mathrm{new}}(a)=y(a)-y(a_0)+p(a_0)
\end{equation}
до clipping. Ошибка baseline остаётся общей ошибкой альтернатив; clipping дополнительно искажает ранжирование. Absolute-пробник в этих условиях получает $p_{\mathrm{new}}(a)=y(a)$. Это объясняет возможное ограничение конкретной формулы, но oracle-head использует запрещённые методу query outcomes и не является обучаемым достижением.

\subsection{Результат дополнения M и absolute-абляции}
На V1 original-absolute, completed-contrast и completed-absolute дали одинаковые $+4{,}5524$ п.п. На V2 все три дали $-0{,}0651$ вместо $-0{,}9579$ у исходного contrast, но не улучшили ни одной сцены. Новая ridge на completed M дала $-0{,}1654/-0{,}8742$ п.п. на V1/V2~\cite{completion}.

Таким образом, дополнение labels и смена contrast на absolute уменьшили отдельную потерю, но не создали устойчиво полезного переноса. Подбор bandwidth и коэффициентов на этих query outcomes был остановлен.

\section{Прямая регрессия полезности действия}
\subsection{Проверенная идея}
Вместо переноса восьми residuals модель стала предсказывать
\begin{equation}
 \widehat{\Delta Q}(s,a)\approx Q(s,a)-Q(s,a_0).
\end{equation}
\code{direct_no_memory} использует scene, baseline и относительные action features. \code{direct_memory} добавляет геометрически взвешенные прошлые $\Delta Q$, поддержку и pooled-оценки M-доноров. Обе модели получают одинаковые полные top16 targets M/F.

Baseline имеет нулевой вектор относительных признаков и прогноз ноль. Ridge-регрессия не использует intercept, нормировка обучается только внутри train-fold; $\lambda\in\{1,10,100,1000\}$ выбирается leave-one-F-session-out. Действие допускается при положительном predicted gain. Это utility-модель, не calibrated safety gate.

\subsection{Результаты до расширения данных}
\par
\begin{table}[htbp]\centering\small
\caption{Scalar pair utility на исходных M48/F32. H --- число relative safety-harm событий.}
\begin{tabular}{@{}lrrrr@{}}\toprule
Метод & Gain V1 & H V1 & Gain V2 & H V2\\\midrule
Direct memory & 0,1213 & 1 & 0,0909 & 0\\
Direct no-memory & -0,0948 & 2 & -0,3197 & 0\\\bottomrule
\end{tabular}
\end{table}
\par
Memory дала небольшой положительный gain на обеих когортах и превзошла этот контроль. Но coverage был $100\%$: все 24 и все 16 решений менялись. На V1 оставался safety-harm; малый gain при почти безусловном вмешательстве не удовлетворяет исходной цели полезного управляемого recall~\cite{pair}.

\section{Разметочная поддержка safety и расширение M/F}
\subsection{Почему нулевой harm не был достаточным}
Аудит исходного top16 без baseline показал слабую и неоднородную поддержку safety-degradation labels. В F среди 480 альтернатив были три ухудшения NC и 14 TTC. В M среди 720 альтернатив не было NC-ухудшений, было пять DAC и одно TTC. У DDC/TLC не было ухудшений относительно baseline в этом top16-срезе, хотя абсолютные нарушения встречались.

В random tail появлялись дополнительные ошибки, но tail ranks $>16$ не совпадают с deployment pool. Поэтому их нельзя использовать как доказательство калибровки риска в top16. Counts относятся к коррелированным альтернативам внутри сцен, а не к числу независимых наблюдений.

\subsection{Как собрали дополнительные примеры}
Из уже имеющихся архивов не нашлось требуемых четырёх новых подходящих сессий без image transfer. Пользователь разрешил до $8{,}5$ ГБ; выбран shard 174 размером $8{,}4276$ ГБ. Подготовлены 64 screening-сцены, а до открытия исходов зафиксированы 32 label-сцены: 16 uniform и 16 challenge, по две новые M- и F-сессии.

Challenge-отбор использовал доступные прогнозы для поиска потенциальных repair/degradation случаев; он не гарантировал наличия измеренных safety-событий. Разметка завершилась $32/32$ сценами, $512$ top16 candidate evaluations и отдельными raw/filtered вызовами. Uniform/challenge результаты сообщались раздельно.

В новых challenge F-альтернативах обнаружены пять NC- и восемь TTC-ухудшений, в challenge M --- 16 DAC-, одно NC- и одно TTC-ухудшение. В uniform M встретились по одному NC/TTC; в uniform F не было ухудшений пяти учитываемых компонентов~\cite{newmf}. Поддержка стала богаче, но оставалась редкой и сосредоточенной в отдельных сессиях. Новые четыре сессии стали M/F, а не независимым тестом.

\subsection{Data-only refit scalar utility}
Замороженная scalar-формула была переобучена на объединённых M64/F48, без смены цели на этой стадии. Получено~\cite{expanded}:
\par
\begin{table}[htbp]\centering\small
\caption{Прямая полезность после M/F-расширения.}
\begin{tabular}{@{}lrrrr@{}}\toprule
Метод & Gain V1 & H V1 & Gain V2 & H V2\\\midrule
Direct memory & -7,9569 & 3 & -8,5264 & 3\\
Direct no-memory & -7,9569 & 3 & -8,5555 & 3\\\bottomrule
\end{tabular}
\end{table}
\par
На V1 методы получили одинаковый агрегат, на V2 различие мало относительно общей деградации. Primary стал baseline. Это сильный отрицательный результат для данной scalar-постановки. Он не доказывает, что больше данных вредно вообще: одновременно изменились состав и распределение M/F, а новый challenge-срез отличается от uniform. Но ожидание, что одних дополнительных safety labels будет достаточно, не подтвердилось.

\section{Ограниченная компонентная регрессия}
\subsection{Цели и модель}
Следующая идея --- отдельно прогнозировать абсолютный PDM и восемь компонентов:
\begin{equation}
 \mathbf y=(Q,NC,DAC,DDC,TLC,EP,TTC,LK,HC).
\end{equation}
Метки обрезаются в $[0{,}001,0{,}999]$, преобразуются logit-функцией, затем обучается multi-output ridge с непенализируемым intercept и sigmoid выхода. Минимизируется squared error в logit-пространстве, а не Bernoulli likelihood. Полученные числа ограничены $[0,1]$, но не объявляются калиброванными вероятностями или confidence bounds.

Признаки содержат full normalized visual embedding, ego history, абсолютные geometry/heads, context--action interactions и gap model score относительно baseline. Recall добавляет девять прошлых исходов, их baseline contrast, donor disagreement и абсолютную поддержку. Важное отличие от scalar модели: baseline имеет собственный абсолютный прогноз.

Четыре family --- no-memory, real memory, random-compatible и shuffled outcomes --- обучаются отдельно на одинаковых истинных targets. Shuffle переставляет целые девятимерные outcome vectors внутри command-групп; истинные supervised targets не перемешиваются. Это один фиксированный random/shuffle seed, а не оценка seed-устойчивости.

\subsection{Правила принятия решения}
\begin{description}
\item[Raw.] Максимизировать predicted PDM, ties разрешать в пользу baseline.
\item[Component.] Среди действий с predicted safety-компонентами не хуже baseline выбрать максимальный predicted PDM.
\item[Guarded.] Дополнительно потребовать неухудшения штатных teacher safety-heads относительно baseline.
\end{description}
Отсутствие совместимого M-донора ведёт к baseline fallback. Допуск $10^{-8}$ --- численный, не вероятностный safety margin. Raw/component/guarded показывают цену инженерных фильтров; одинаковый predicted safety не гарантирует одинакового фактического safety.

\subsection{F-CV и выбор primary}
Во всех folds нормировки и intercept обучаются только на train. Регуляризация выбирается из той же сетки по guarded F-CV: сначала предпочтение нулевому observed harm, затем uniform macro-session gain, далее tie-breaks. Challenge harm учитывается, но challenge gain не подменяет uniform quality.

Primary выбирается только среди memory-guarded/no-memory-guarded при $H_F=0$ и положительном uniform gain. При отсутствии подходящего варианта выбирается baseline, в точной ничьей --- no-memory. Это правило было зафиксировано до оценки V и не должно заменяться постфактум по понравившемуся V-результату.

\subsection{Полученный результат}
Memory raw дал $+1{,}0918/-6{,}3297$ п.п. на V1/V2 с одним harm в каждой когорте. Memory component дал $-0{,}2339/0$, memory guarded $-0{,}1367/0$. Guarded вмешался в двух V1 сценах, обе ухудшили PDM; на V2 не вмешивался. No-memory raw дал $+1{,}1813/-6{,}0928$, random raw $+1{,}1144/-5{,}6101$~\cite{component}. Primary остался baseline. Ограниченные выходы и pointwise gates уменьшили вмешательства, но не обеспечили полезную память.

\section{Последняя абляция: candidate-conditioned decoder features}
\subsection{Гипотеза}
Проверялась возможность, что усреднённый scene embedding и handcrafted geometry теряют важную информацию о конкретном действии. Из последнего refinement decoder layer извлекался latent $z_a$ каждого кандидата. Ко всем четырём семействам одинаково добавлялись:
\begin{equation}
 [\,z_a,\ z_a-z_{a_0}\,].
\end{equation}
На сервере tensor имел форму $256\times256$ на сцену; добавляются 512 признаков на действие. Encoder не дообучается, PCA не вводится. Labels, ridge loss, CV и политики raw/component/guarded сохраняются. Retrieval остаётся на прежнем scene embedding/geometry: эта абляция не меняет одновременно и представление для поиска, и оценку действия.

Экспорт включает ровно 152 сцены: 80 исходных M/F + 24 V1 + 32 новых M/F + 16 V2. Исходные C/T и неразмеченные 32 screening-сцены расширения не включены. Кодовые тесты сравнивали AST ключевых learner-функций с прежней компонентной версией.

\subsection{Первоначальные проблемы паритета}
Обычный повторный forward мог менять \code{candidate_ids}. Исходный sidecar сохранил девять сцен и остановился. Диагностика в отдельном запуске иногда проходила, однако последующий полный экспорт снова падал. Это не было доказательством стабильного естественного coarse top-k.

Введён отдельный conditional replay: перед refinement подставляются сохранённые 256 IDs в прежнем порядке, их vocabulary trajectories, соответствующие полные coarse states и logits. Естественный coarse drift записывается отдельно. Поведение всего естественного candidate-generation pipeline этим не подтверждается; проверяется представление на фиксированном пуле.

Replay затем сохранил 25 сцен и остановился на IMI numerical parity. Read-only probe показал небольшие расхождения даже между текущими проходами. Увеличивать число попыток до случайного прохождения было бы неверно; вместо этого введён новый фиксированный численный режим.

\section{Численная воспроизводимость и пересмотр критериев приёмки}
\subsection{Детерминированный режим}
Перед Python фиксируются seed, один CUDA device и \code{CUBLAS_WORKSPACE_CONFIG=:4096:8}. До CUDA initialization включается \code{torch.use_deterministic_algorithms(True, warn_only=False)}, отключаются cuDNN benchmark/TF32, включается deterministic cuDNN и math-only PyTorch SDP attention. Фиксируются версии Torch/CUDA/cuDNN, GPU и observed flags. Подход основан на рекомендациях PyTorch 2.0.1~\cite{torch}.

Два текущих forward --- без и с decoder-hook --- должны совпасть точно по всем сохраняемым heads, scores, IDs, action и trajectory. Между ними seed не сбрасывается. Отдельно heads повторно применяются к captured latent. Эти проверки сильнее простого allclose между текущими проходами.

Детерминированные текущие outputs всё же немного отличались от ранее сохранённых outputs другого вычислительного режима. Например, на \code{dbb6af7aa1415da4} один DAC-logit и два IMI-logit вышли за старый допуск при точном current repeat и неизменном baseline. Это различие численных реализаций, а не установленная ошибка hook.

\subsection{Промежуточный numeric contract и полный аудит}
Отдельный контракт зафиксировал raw-logit drift $\le5\cdot10^{-5}$, metric-sigmoid drift $\le10^{-5}$, IMI log-softmax drift $\le10^{-4}$, прежний допуск scores и точные baseline/top16. Он остановился на IMI drift $5{,}2452\cdot10^{-5}$. Вместо очередного повышения границы выполнен аудит всех 152 сцен без приёмки features и без fit.

\par
\begin{table}[htbp]\centering\small
\caption{Полный численный аудит. Максимумы рассчитаны по 152 сценам.}
\begin{tabularx}{\linewidth}{@{}Xr@{}}\toprule
Проверка & Результат\\\midrule
Проверено сцен & 152/152\\
Полностью прошли промежуточный contract & 150/152\\
Current repeat failures & 0\\
Baseline / top16 / vocabulary changes & 0 / 0 / 0\\
Decoder/head alignment failures & 0\\
Сцены вне старого strict tolerance & 63\\
Максимальный score drift & $8{,}5831\cdot10^{-6}$\\
Максимальный metric sigmoid drift & $3{,}0534\cdot10^{-6}$\\
Максимальный IMI-logit drift & $1{,}0109\cdot10^{-4}$\\
Максимальный IMI log-softmax drift & $1{,}1568\cdot10^{-4}$\\\bottomrule
\end{tabularx}
\end{table}
\par
Два нарушения intermediate contract: \code{e7832ec3cbcb5fe4} (M) и \code{fedf56bd69af5cec} (F). Они сохранены в аудите, а не удалены или пересчитаны до удачного результата. Счётчик \code{legacy_failed_elements} в field-maxima --- максимум внутри одной сцены, не сумма по всем 152 сценам.

\subsection{Финальный функциональный критерий}
После полного аудита пользователь явно одобрил отдельный режим функциональной эквивалентности. В нём:
\begin{enumerate}
\item остаются точный current repeat, vocabulary, порядок всех 256 IDs, baseline, trajectory, состав top16;
\item score drift сохраняет прежние $\mathrm{atol}=10^{-5}$, $\mathrm{rtol}=10^{-6}$;
\item все восемь metric heads сохраняют пределы raw-logit $5\cdot10^{-5}$ и sigmoid $10^{-5}$;
\item IMI обязан быть конечным и иметь верную форму; raw и normalized drift записываются, но отдельно не порогируются;
\item текущая IMI-head по-прежнему входит в exact-repeat и latent/head alignment проверки.
\end{enumerate}
Исключений по отдельным токенам нет. Критерий \textbf{пересмотрен после аудита}; его нельзя выдавать за заранее объявленный strict parity. Для принятого экспорта связываются hashes полного prior audit и всех 152 записей, включая failures. Старые roots/identity/failure сохранены. Новый root --- \code{decoder_features_functional_v1}.

В learner остаются прежние frozen heads, model scores, top16 и labels; из нового forward используется только latent. Это делает абляцию интерпретируемой, но не доказывает bitwise воспроизводимость естественного старого pipeline или устойчивость на другом hardware. При экспорте каждое functional-условие проверялось заново. Успешный финальный fit/evaluate report означает, что требование полного 152-scene sidecar выполнено.

\subsection{Другие инженерные ошибки, которые были устранены}
Встречались CUDA unavailable/NVML ошибки контейнера при работающем host GPU; их нужно отличать от научного failure метода. Диагностика сохранила начатые артефакты. Ошибки путей/read-only filesystem решены отдельным writable data root.

Отдельно исправлена ошибка \code{check}: вложенные launcher-тесты наследовали root переменные родительского эксперимента и молча падали на assertions. Поэтому отсутствие вывода \code{check \&\& plan \&\& export-pilot} не означало успех пилота. Тесты теперь очищают DRIVEMEM-переменные только в дочернем процессе, печатают stage/строку ошибки, а check явно сообщает об успехе. Добавлена регрессия с намеренно чужими roots/GPU settings. Научные gates и старые экспериментальные файлы этим исправлением не менялись.

\section{Финальные результаты decoder-абляции}
\subsection{Основное сравнение}
Все методы используют один и тот же supervision и добавленные decoder features. Таблица содержит baseline-relative macro-session gain. Полная таблица четырёх семейств и трёх политик приведена в приложении~\ref{app:full}.

\par
\begin{table}[htbp]\centering\small
\caption{Decoder-абляция: gain относительно baseline и observed safety-harm.}
\begin{tabular}{@{}lrrrr@{}}\toprule
Метод & Gain V1 & H V1 & Gain V2 & H V2\\\midrule
\code{memory_guarded} & +4,0653 & 0 & +0,0000 & 0\\
\code{no_memory_guarded} & -0,1582 & 0 & +0,0000 & 0\\
\code{memory_component} & +4,0653 & 0 & +0,0865 & 0\\
\code{no_memory_component} & -0,2134 & 0 & +0,1060 & 0\\
\code{random_component} & +3,8794 & 0 & -0,2640 & 0\\
\code{shuffled_component} & -0,2965 & 0 & -0,0143 & 0\\
\code{memory_raw} & +1,2337 & 1 & -6,4173 & 1\\
\code{no_memory_raw} & +0,9414 & 1 & -7,5758 & 1\\
\bottomrule\end{tabular}\end{table}
\par

\code{memory_guarded} даёт $+4{,}0653$ п.п. на V1, но ноль на V2 с нулевым coverage. \code{memory_component} на V2 даёт $+0{,}0865$ п.п., меньше no-memory-component ($+0{,}1060$). Raw memory на V2 теряет $6{,}4173$ п.п. и имеет один safety-harm. Random-component на V1 даёт $+3{,}8794$ п.п., близко к memory-component; это дополнительно ослабляет вывод о специфической пользе similarity retrieval~\cite{decoder}.

Парный guarded memory-minus-no-memory выигрыш на V1 равен $4{,}2236$ п.п. с session CI $[0;12{,}5551]$; на V2 он нулевой. CI включает ноль, а две маленькие development-когорты уже участвовали в анализе предыдущих гипотез.

\subsection{Почему primary остался baseline}
Из frozen fit~\cite{finaldiagnostic}:
\par
\begin{table}[htbp]\centering\small
\caption{Результат выбора по F-CV для последней абляции.}
\begin{tabular}{@{}lrrrr@{}}\toprule
Семейство & $\lambda$ & Uniform gain, п.п. & H & Вмешательства\\\midrule
No-memory & 1000 & 0 & 0 & 3\\
Memory & 100 & 0 & 0 & 2\\
Random & 1000 & 0 & 0 & 2\\
Shuffled & 1000 & 0 & 0 & 4\\\bottomrule
\end{tabular}
\end{table}
\par
У обоих допустимых претендентов на primary отсутствует положительный uniform F-CV gain. Условие $H_F=0$ выполнено, условие положительной полезности --- нет. Правило выбора корректно оставило baseline. Заменить primary на memory-guarded после просмотра положительного V1 было бы post-hoc selection по development.

\subsection{Разбор положительного и отрицательного вмешательства V1}
В \code{memory_guarded} всего два вмешательства, оба в одной сессии \code{2021.10.11.07.47.13_veh-50}:
\begin{description}
\item[Положительное.] Token \code{f27a618a5c7458ee}, candidate 6143. Local PDM $0\rightarrow1$, gain $+100$ п.п. DAC меняется $0\rightarrow1$. NC, DDC, TLC, TTC, progress, lane keeping и comfort не меняются. Следовательно, это исправление drivable-area compliance по метрике, а не свидетельство исправленного столкновения.
\item[Отрицательное.] Token \code{dbfedc756de85252}, candidate 249. PDM $0{,}91945\rightarrow0{,}895131$, gain $-2{,}43191$ п.п. Progress уменьшается на $0{,}068093$, safety-компоненты не меняются; поэтому $H=0$, но quality degradation есть.
\end{description}
Общий выигрыш выводится непосредственно:
\begin{equation}
 \frac{100-2{,}43191}{24}=4{,}06534\ \text{п.п.}
\end{equation}
Три session gains: $0$, $0$, $12{,}196011$ п.п. После удаления последней сессии gain становится нулевым; при удалении любой из двух других --- $6{,}098006$ п.п. Это не переобучение и не новый тест, а проверка концентрации эффекта.

Тот же token \code{f27a618a5c7458ee} уже давал $+100$ п.п. в action-transfer B3 с другим выбранным candidate ID. Разные модели повторно находят известное исправление. Такие повторения не являются независимыми успешными случаями. Component-политика добавила ещё три смены действий с нулевым gain, поэтому её средний результат V1 совпал с guarded.

\subsection{Что произошло на V2}
Guarded memory не меняет ни одного из 16 решений. Component memory улучшает только progress в двух сценах:
\par
\begin{table}[htbp]\centering\small
\caption{Положительные component-вмешательства на V2. Safety-компоненты не изменились.}
\begin{tabular}{@{}lrrr@{}}\toprule
Token & Candidate & Gain, п.п. & $\Delta EP$\\\midrule
\code{c2f8efa8358050de} & 6536 & 0,057670 & 0,001615\\
\code{3f526a6123c157e4} & 7074 & 1,326653 & 0,037146\\\bottomrule
\end{tabular}
\end{table}
\par
Их сумма, делённая на 16, даёт $0{,}086520$ п.п. Из четырёх session gains два нулевые, остальные $0{,}014417$ и $0{,}331663$ п.п. После исключения последней сессии остаётся лишь $0{,}004806$ п.п. на остальных трёх. Даже этот небольшой перенос не превосходит no-memory-component на V2.

\section{Сводная оценка проверенных идей}
\begin{longtable}{@{}p{0.27\linewidth}p{0.33\linewidth}p{0.33\linewidth}@{}}
\caption{Что проверяли и что позволяет заключить результат.}\label{tab:hypotheses}\\
\toprule Гипотеза & Проверка & Итог\\\midrule\endfirsthead
\toprule Гипотеза & Проверка & Итог\\\midrule\endhead
Кандидаты содержат резерв & Oracle top256/top16 & Подтверждено на диагностических наборах; не результат памяти.\\
Похожесть переносит head residuals & Similarity/random/threshold & Similarity gain нулевой в первом пилоте; полезность из близости не следует.\\
Нужны локальные action kernels/contrast & B0--B3 & Большой V1-плюс, концентрация в одной сессии, отрицательный V2.\\
Gate --- главный bottleneck & Gate-oracle/family-oracle & Удаление вредных решений недостаточно: бедны сами предложения семейства.\\
Виновата неполная top16-разметка M & Дополнение 369 labels & Уменьшена одна потеря, устойчивые улучшения не появились.\\
Contrast ограничивает score & Ceiling и oracle-head probes & Ограничение конкретного contrast target обнаружено, но learned utility этим не доказана.\\
Прямой $\Delta$PDM проще учить & Scalar pair utility & Небольшие плюсы с почти полным вмешательством; harm и неустойчивость.\\
Дополнительные safety labels решат проблему & Новые M/F и data-only refit & Поддержка выросла, scalar модели сильно деградировали.\\
Компоненты и bounded outputs дадут полезный отказ & Девять targets, три gates & Вред уменьшается, но часто вместе с пользой; primary baseline.\\
Нужна candidate-conditioned representation & Decoder feature ablation & Локальный известный repair, без устойчивого преимущества на V2/F-CV.\\
Проблема только в hook/GPU недетерминизме & Exact-repeat и полный audit & Техническая воспроизводимость достигнута на проверенных сценах; научный utility failure сохранился.\\
\bottomrule
\end{longtable}

\section{Положительные результаты, отрицательные результаты и ограничения}
\subsection{Что является самостоятельным полезным результатом}
Получен работающий стенд с контролем происхождения данных и парным сравнением действий, выявлены опасные места human-filter snapshots и псевдорепликации по log-segments, реализован строгий M-only recall и сильные no-memory контроли. Существование action-oracle резерва подтверждено на нескольких наборах. Найден воспроизводимый local repair DAC-сцены. Временные и денежные ограничения переведены в явные budgets, manifests и ограниченные планы загрузки.

Полезным является и отрицательный методический результат: одинаковый доступ к supervised labels, случайные доноры, shuffled outcomes и разбор одной сессии существенно меняют интерпретацию среднего выигрыша. Без этих контролей $+4$ п.п. на V1 можно было бы ошибочно представить как устойчивое улучшение DriveMem.

\subsection{Что не следует объявлять доказанным}
\begin{enumerate}
\item Oracle gain не является достижимым learned gain и не оценивает причинный эффект на реальной дороге.
\item Нулевой $H$ не является calibrated risk control, особенно при нулевом coverage или редких событиях.
\item Sigmoid-выходы ridge и teacher не являются автоматически вероятностями вреда.
\item Уже просмотренные V1/V2 не являются fresh holdout. C/T независимы от обучения нового модуля по роли, но их независимость от pretraining базового checkpoint не доказана.
\item Enriched frozen-bank pilot не доказывает пользу executed-only или continually growing memory.
\item Functional export на сохранённом пуле не доказывает прежний strict parity всего естественного pipeline.
\item Смена численного критерия после аудита должна быть раскрыта; это инженерное решение, не новый статистический сертификат.
\item По трём и четырём сессиям нельзя делать сильные выводы о редких рисках или устойчивости на популяции.
\end{enumerate}

\subsection{Почему отрицательные результаты не опровергают всю исходную идею}
Проверялись небольшие fixed banks, линейные модели, конкретные features и геометрические kernels. Не исследованы полноценная динамика взаимодействий, другая архитектура utility predictor, систематический рост банка и независимые новые географии. Поэтому корректный вывод относится к реализованным вариантам и их данным. При этом повторяющееся отсутствие преимущества над no-memory и концентрация в одном известном случае являются основанием остановить текущий цикл подбора, а не продолжать его бесконечно на тех же V.

\section{Сопоставление с шестимесячным планом статьи}
\par
\begin{table}[htbp]\centering\small
\caption{Статус ключевых целей исходного плана.}
\begin{tabularx}{\linewidth}{@{}p{0.47\linewidth}X@{}}\toprule
Цель & Статус на 07.11.2026\\\midrule
Frozen planner и candidate reranking & Реализовано; v2 reproduction и oracle diagnostics\\
Парные компонентные effects & Реализованы local targets и отчёты raw/filtered\\
Положительные и вредные memory interventions & Наблюдаются; устойчивой пользы не установлено\\
Сопоставимые no-memory/random controls & Реализованы; ограничивают положительные claims\\
Уверенность, lower bounds и calibration & Не реализованы как подтверждённый риск-контроль\\
Полезный selective recall & Инженерные gates есть; полезность и risk guarantee не подтверждены\\
Рост памяти во времени & Не проверен; банк фиксирован\\
Executed-only против enriched & Executed-only механизм предусмотрен, сравнительный результат не получен\\
Новые карты, OOD и второй planner & Не проверены\\
Полный EPDMS с памятью, closed loop & Не измерены\\\bottomrule
\end{tabularx}
\end{table}
\par

Текст о стенде, oracle motivation, отрицательном переносе, контролях и ограничениях можно писать уже сейчас. Исходное название с \emph{Risk-Controlled} и \emph{Continual} описывает исследовательскую цель, а не полученный метод. Отрицательный пилот сам по себе не гарантирует достаточную новизну для целевой конференции; потребуется переоценка постановки, литературы и вклада отдельно от этого отчёта.

\section{Рекомендуемое завершение текущей ветки и дальнейшие решения}
\subsection{Что зафиксировать сейчас}
Сохранить все frozen roots, manifests, fit, decision sets и per-scene отчёты; оставить primary baseline. Текущую decoder-абляцию завершить без изменения gates по V, дополнительного перебора thresholds и выбора удачного seed. Зафиксировать DAC-repair как case study вместе с отрицательным progress-вмешательством и случаями нулевого coverage.

\subsection{Если исследование продолжается}
Следующее направление должно заранее задавать новую проверяемую причину неудачи: ограниченность causal/interaction context, различие action support и outcome transferability, либо отсутствие достаточно разнообразных training effects. Изменение targets, представления или дизайна данных должно иметь такой же сильный no-memory контроль и новый, не использованный при выборе этой гипотезы evaluation протокол.

Калибровка риска имеет смысл после появления ненулевой полезности и достаточного числа независимых вмешательств. Она не может создать utility там, где метод почти всегда отказывается. Открывать C/T для дальнейшего tuning вместо отдельного исследования нельзя. При положительном новом development-сигнале нужны заранее зафиксированный метод, независимая проверка, затем full EPDMS на последовательностях, executed-only/growing-bank и второй planner.

\resultbox{\textbf{Итоговая научная формулировка.} В изученном малом NAVSIM-пилоте у DriveSuprim есть заметный резерв качества внутри готовых кандидатов. Реализованные эпизодические memory-rerankers могут исправлять отдельные ошибки выбора, но устойчивое преимущество над сопоставимыми no-memory моделями не установлено. Компонентные ограничения уменьшают вмешательства и вред, однако не дают подтверждённого полезного risk-controlled recall.}

\clearpage
\appendix
\section{Подробные численные таблицы}\label{app:full}
Все gains в приложении --- п.п. local PDM относительно baseline своей когорты. H --- число сцен с падением хотя бы одного из пяти учитываемых safety-компонентов. I --- число вмешательств. Таблицы сгенерированы из переданных JSON, включённых в \code{evidence_extracted.json}; округление выполнено только для отображения.

\subsection{Action transfer B0--B3 и контроли}
{\footnotesize\setlength{\tabcolsep}{4pt}
\begin{longtable}{@{}lrrrrrr@{}}
\toprule & \multicolumn{3}{c}{V1} & \multicolumn{3}{c}{V2}\\
Метод & Gain & H & I & Gain & H & I\\\midrule\endhead
\code{B0} & +0,0000 & 0 & 2 & +0,0000 & 0 & 1\\
\code{B1} & +4,5524 & 0 & 7 & +0,0000 & 0 & 1\\
\code{B2} & +0,0000 & 0 & 0 & +0,0000 & 0 & 0\\
\code{B3} & +4,2356 & 0 & 4 & -0,9579 & 0 & 4\\
\code{no_memory_ridge_full} & +1,3234 & 0 & 17 & -0,8742 & 0 & 11\\
\code{random_0} & +4,0506 & 0 & 5 & -0,8929 & 0 & 2\\
\code{random_1} & -0,0107 & 0 & 6 & +0,0805 & 0 & 3\\
\code{random_2} & +4,0893 & 0 & 4 & -0,9579 & 0 & 2\\
\code{scene_only} & +0,0000 & 0 & 2 & +0,0000 & 0 & 1\\
\bottomrule\end{longtable}}

\subsection{Дополнение M top16 и абсолютный перенос}
{\footnotesize\setlength{\tabcolsep}{4pt}
\begin{longtable}{@{}lrrrrrr@{}}
\toprule & \multicolumn{3}{c}{V1} & \multicolumn{3}{c}{V2}\\
Метод & Gain & H & I & Gain & H & I\\\midrule\endhead
\code{original_M_contrast} & +4,2356 & 0 & 4 & -0,9579 & 0 & 4\\
\code{original_M_absolute} & +4,5524 & 0 & 5 & -0,0651 & 0 & 3\\
\code{completed_M_contrast} & +4,5524 & 0 & 4 & -0,0651 & 0 & 4\\
\code{completed_M_absolute} & +4,5524 & 0 & 6 & -0,0651 & 0 & 4\\
\code{ridge_original} & +1,3234 & 0 & 17 & -0,8742 & 0 & 11\\
\code{ridge_completed_M} & -0,1654 & 0 & 16 & -0,8742 & 0 & 12\\
\bottomrule\end{longtable}}

\clearpage
\subsection{Компонентная модель до decoder features}
{\footnotesize\setlength{\tabcolsep}{4pt}
\begin{longtable}{@{}lrrrrrr@{}}
\toprule & \multicolumn{3}{c}{V1} & \multicolumn{3}{c}{V2}\\
Метод & Gain & H & I & Gain & H & I\\\midrule\endhead
\code{baseline} & +0,0000 & 0 & 0 & +0,0000 & 0 & 0\\
\code{memory_component} & -0,2339 & 0 & 4 & +0,0000 & 0 & 2\\
\code{memory_guarded} & -0,1367 & 0 & 2 & +0,0000 & 0 & 0\\
\code{memory_raw} & +1,0918 & 1 & 21 & -6,3297 & 1 & 15\\
\code{no_memory_component} & -2,8401 & 2 & 8 & -0,0662 & 0 & 2\\
\code{no_memory_guarded} & -0,1399 & 0 & 3 & -0,2583 & 0 & 1\\
\code{no_memory_raw} & +1,1813 & 1 & 22 & -6,0928 & 1 & 15\\
\code{random_component} & -2,8144 & 1 & 4 & -0,0662 & 0 & 2\\
\code{random_guarded} & -0,0386 & 0 & 1 & +0,0000 & 0 & 0\\
\code{random_raw} & +1,1144 & 1 & 20 & -5,6101 & 1 & 16\\
\code{shuffled_component} & -0,1527 & 0 & 3 & -0,1805 & 0 & 1\\
\code{shuffled_guarded} & -0,1197 & 0 & 1 & -0,2583 & 0 & 1\\
\code{shuffled_raw} & +0,6854 & 1 & 19 & -5,6777 & 1 & 11\\
\bottomrule\end{longtable}}

\subsection{Decoder features: все семейства и политики}
{\footnotesize\setlength{\tabcolsep}{4pt}
\begin{longtable}{@{}lrrrrrr@{}}
\toprule & \multicolumn{3}{c}{V1} & \multicolumn{3}{c}{V2}\\
Метод & Gain & H & I & Gain & H & I\\\midrule\endhead
\code{baseline} & +0,0000 & 0 & 0 & +0,0000 & 0 & 0\\
\code{memory_component} & +4,0653 & 0 & 5 & +0,0865 & 0 & 2\\
\code{memory_guarded} & +4,0653 & 0 & 2 & +0,0000 & 0 & 0\\
\code{memory_raw} & +1,2337 & 1 & 20 & -6,4173 & 1 & 14\\
\code{no_memory_component} & -0,2134 & 0 & 4 & +0,1060 & 0 & 4\\
\code{no_memory_guarded} & -0,1582 & 0 & 2 & +0,0000 & 0 & 1\\
\code{no_memory_raw} & +0,9414 & 1 & 22 & -7,5758 & 1 & 14\\
\code{random_component} & +3,8794 & 0 & 6 & -0,2640 & 0 & 4\\
\code{random_guarded} & -0,1582 & 0 & 2 & +0,0000 & 0 & 1\\
\code{random_raw} & +0,8969 & 1 & 23 & -7,5095 & 1 & 16\\
\code{shuffled_component} & -0,2965 & 0 & 7 & -0,0143 & 0 & 3\\
\code{shuffled_guarded} & -0,1582 & 0 & 2 & +0,0000 & 0 & 1\\
\code{shuffled_raw} & +0,9751 & 1 & 22 & -7,4252 & 1 & 14\\
\bottomrule\end{longtable}}

\clearpage
\section{Артефакты и воспроизводимость}
\subsection{Основные каталоги серверных экспериментов}
Общий префикс --- \path{/workspace/exp_v2/drivemem_memory/}. Относительные имена:
\begin{longtable}{@{}p{0.46\linewidth}p{0.48\linewidth}@{}}\toprule
Каталог & Содержание\\\midrule\endhead
\path{pilot_v1} & Полная инвентаризация и первоначальный 352-scene split/cache\\
\path{budget20gb_v1} & Исходные M/F/V1/C/T, features и cache\\
\path{validation_cached4_v1} & V2, 16 сцен и четыре сессии\\
\path{mf_new4_shard174_v1} & Новые M/F metadata, screening, labels и support report\\
\path{experiments/action_transfer_v2} & B0--B3 и сильный ridge контроль\\
\path{experiments/memory_top16_v1} & Дополнение M labels и absolute/contrast сравнение\\
\path{experiments/pair_utility_v1} & Исходная scalar pair utility\\
\path{experiments/pair_utility_expanded_v1} & Data-only refit на M64/F48\\
\path{experiments/component_utility_v1} & Девять абсолютных targets, три политики\\
\path{decoder_numeric_audit_v1} & Полный 152-scene численный аудит\\
\path{decoder_features_functional_v1} & Принятые функциональным критерием decoder sidecars\\
\path{experiments/component_decoder_functional_v1} & Финальный decoder-feature эксперимент\\\bottomrule
\end{longtable}

Необходимо сохранять и отклонённые strict/replay/deterministic/numeric roots, чтобы история приёмки не исчезла. Имена \code{validation} и \code{test} сами по себе не гарантируют независимость данных; важна история фактического использования.

\subsection{Основные реализованные кодовые пакеты}
\begin{itemize}
\item \path{navsim/agents/drivemem.py}: базовые метрики, retrieval и residual memory.
\item \path{navsim/agents/drivemem_action_transfer.py}: action geometry, B0--B3 и ridge features.
\item \path{navsim/agents/drivemem_pair_utility.py}: прямая относительная полезность.
\item \path{navsim/agents/drivemem_component_utility.py}:\\
девять targets, bounded regression, component/guarded policies.
\item \path{navsim/agents/drivemem_decoder_utility.py}: общий decoder-feature интерфейс и проверка sidecar provenance.
\item \path{navsim/agents/drivemem_determinism.py}, \path{drivemem_numeric_contract.py}, \path{drivemem_functional_contract.py}: численный режим, разные версии критериев приёмки и связь с аудитом. Последние два имени находятся в том же каталоге agents.
\item \path{navsim/planning/script/run_drivemem_*}: inventory/cache/export/labels/fit/evaluate и отдельные export/audit протоколы.
\item \path{scripts/drivesuprim/evaluation/}: launchers, tests, per-scene/support/family/safety/numeric diagnostics.
\item \path{download/}: metadata preparation, archive plans, bounded image acquisition и новые M/F splits.
\end{itemize}

\subsection{Чтение финального результата на сервере}
\begin{verbatim}
docker exec egorov_drivesuprim cat \
  /workspace/exp_v2/drivemem_memory/experiments/\
component_decoder_functional_v1/report.json
\end{verbatim}
Обратная косая черта в конце строки соединяет части пути при исполнении в Bash. Удобный launcher:
\begin{verbatim}
docker exec -w /workspace/DriveSuprim egorov_drivesuprim \
  bash scripts/drivesuprim/evaluation/\
run_drivemem_decoder_functional.sh report
\end{verbatim}
При наличии frozen артефактов эти команды читают результат; повторное обучение ради написания статьи не требуется. Дополнительный разбор выводится \path{scripts/drivesuprim/evaluation/summarize_drivemem_decoder_result.py}. Нельзя запускать новый writer в существующий frozen root после изменения кода.

\section{Источники, полнота данных и контроль происхождения}
Отчёт составлен по состоянию доступного локального репозитория и чата на 7 ноября 2026 года. Git HEAD при подготовке: \code{9d3ef022445a0e7cc2de9298a45b3e41dce265d2}. Это идентификатор локального среза, а не доказательство, что каждый ранний серверный запуск использовал именно этот commit: для этого служат его собственные identities.

В \code{source_manifest.json} сохранены SHA256 исходного PDF, выбранных документов/модулей и присланных вложений. В \code{evidence_extracted.json} сохранены распознанные верхнеуровневые JSON-объекты серверных выводов. Это не копия полного серверного experiment directory. Часть ранних oracle и полного численного audit результатов была передана непосредственно текстом чата; они явно относятся к сообщённым пользователем результатам.

Точные промежуточные размеры/CI и дополнительные family-capacity числа сообщаются только там, где они доступны в логах или документации. Не восстановлены отсутствующие per-candidate distributions, фактические wall times всех этапов, стоимость всей сети, GPU peak memory и полный session-bootstrap для раннего 32-log oracle. Эти величины не следует выдумывать или выводить из одного агрегата.

\clearpage
\begin{thebibliography}{99}
\raggedright
\bibitem{plan} Исходный документ пользователя. \emph{DriveMem: шестимесячный план подготовки статьи для NeurIPS 2027}. Версия 15.08.2026. \code{DriveMem_план_статьи_6_месяцев.pdf}. Использован как план, не как подтверждение выполненных экспериментов.
\bibitem{repro} Документация проекта. \path{docs/navsim_v1_reproduction_findings.md}; \path{docs/drivemem_revised_research_plan.md}. Числа воспроизведения, зафиксированные в проекте; актуальная литература не переоценивалась в рамках этого отчёта.
\bibitem{initial} Серверный вывод исходного budget-пилота, вложение \code{bed09855-c39b-4096-8a06-26371c2bd16b}. Дополнительная диагностика: \code{1987459a-8056-47c9-8fbe-2d3bb1550133}.
\bibitem{action} Action-transfer report, вложение \code{e024e701-9adb-4b96-93b2-bf5c27be5d00}. Аудит крупнейшего эффекта: \code{d5f3f03f-65ce-493a-ba25-8ca587fd92b1}.
\bibitem{cached} Cached4 report, вложение \code{5ad39327-57a3-4e42-8728-40d1345e5f13}.
\bibitem{support} Action-support audit, вложение \code{23e027fa-75c3-4ebc-af3b-b088fbe72e63}; donor-family audit \code{09d892eb-790a-4713-9731-343e38e9116d}.
\bibitem{ceiling} Score-ceiling и oracle-head diagnostics, вложение \code{042c1568-25bc-4c9a-9ef7-917a92f72f7e}.
\bibitem{completion} M-top16 completion report, вложение \code{36692838-3025-4e35-a423-c34b4c64644a}.
\bibitem{capacity} Completed-family capacity audit, вложение \code{92f65e97-eea0-4b50-9ed3-e2dc5618733d}.
\bibitem{pair} Pair-utility report и fit diagnostics, вложения \code{b6b4448e-182e-4f93-9431-6fef3d87c57f}, \code{8a88a7e5-8aa6-47ce-8c4e-e2dfd973962b}. Safety audit: \code{9386cac6-bff0-4073-ae50-f922f3345f58}.
\bibitem{newmf} Новый M/F support report, вложение \code{09bea6aa-6874-40ef-b20b-9bd4a2f8c372}; протокол \path{docs/drivemem_new_mf_run.md}.
\bibitem{expanded} Expanded-utility report, вложение \code{c6fc3240-4b0a-49d2-be7b-7fc9321238a1}.
\bibitem{component} Component-utility report, вложение \code{b17abf93-c7d2-4cae-9de7-1cb86f5c9a54}.
\bibitem{decoder} Final decoder-functional report, вложение \code{f38b7c94-f0d8-428b-b4a0-4067a5f4f722}.
\bibitem{finaldiagnostic} Frozen F-CV selection, interventions и session sensitivity, вложение \code{f4a39b30-4dfb-4312-86e1-738d8ec028ab}.
\bibitem{torch} PyTorch 2.0.1. \emph{Reproducibility}. \url{https://github.com/pytorch/pytorch/blob/v2.0.1/docs/source/notes/randomness.rst}. Использовано для фиксации численного режима; не источник чисел качества DriveMem.
\end{thebibliography}
\end{document}
